% Figure: superposition of elementary potential flows. Four panels:
% uniform, source, vortex, and uniform+doublet (cylinder).
\begin{figure}[htbp]
\centering
\begin{tikzpicture}[
  sline/.style={draw=engblue, thick},
  body/.style={draw=engblack, fill=enggray!25, thick},
  pt/.style={draw=engred, fill=engred, circle, inner sep=1.3pt},
  label/.style={font=\scriptsize},
  arr/.style={draw=engblue, -{Stealth}, thick},
]
% ===== Panel 1: uniform flow =====
\begin{scope}[xshift=0cm,yshift=0cm]
\draw[draw=enggray, thin] (-1.5,-1.5) rectangle (1.5,1.5);
\foreach \y in {-1.1,-0.55,0,0.55,1.1} {
  \draw[arr] (-1.3,\y) -- (1.3,\y);
}
\node[label, anchor=north] at (0,-1.6) {uniform $U_\infty$};
\end{scope}

% ===== Panel 2: source =====
\begin{scope}[xshift=4cm,yshift=0cm]
\draw[draw=enggray, thin] (-1.5,-1.5) rectangle (1.5,1.5);
\foreach \a in {0,30,...,330} {
  \draw[arr] ({0.25*cos(\a)},{0.25*sin(\a)}) -- ({1.25*cos(\a)},{1.25*sin(\a)});
}
\node[pt] at (0,0) {};
\node[label, anchor=north] at (0,-1.6) {source $m$};
\end{scope}

% ===== Panel 3: point vortex =====
\begin{scope}[xshift=8cm,yshift=0cm]
\draw[draw=enggray, thin] (-1.5,-1.5) rectangle (1.5,1.5);
\foreach \r in {0.45,0.85,1.25} {
  \draw[sline] (0,0) circle (\r);
}
\draw[arr] (1.25,0) -- (1.25,0.4);
\draw[arr] (0.85,0) -- (0.85,0.3);
\node[pt] at (0,0) {};
\node[label, anchor=north] at (0,-1.6) {vortex $\Gamma$};
\end{scope}

% ===== Panel 4: uniform + doublet = cylinder =====
\begin{scope}[xshift=12cm,yshift=0cm]
\draw[draw=enggray, thin] (-1.5,-1.5) rectangle (1.5,1.5);
% streamlines bending around a circle of radius R=0.5: smooth bulged curves.
% Each curve has baseline height \yo, bulging outward by a Gaussian near x=0.
\foreach \yo in {0.55,0.85,1.2} {
  \pgfmathsetmacro{\bulge}{0.35/(\yo)}
  \draw[sline] plot[domain=-1.5:1.5, samples=80, variable=\x]
    ({\x},{ \yo + \bulge*exp(-(\x*\x)/0.45) });
  \draw[sline] plot[domain=-1.5:1.5, samples=80, variable=\x]
    ({\x},{ -\yo - \bulge*exp(-(\x*\x)/0.45) });
}
\draw[body] (0,0) circle (0.5);
\node[label, anchor=north] at (0,-1.6) {uniform $+$ doublet};
\end{scope}

\end{tikzpicture}
\caption[Elementary potential flows and their superposition]{Building
blocks of two-dimensional potential flow. Each is a solution of
Laplace's equation $\nabla^2\phi=0$, so any sum of them is also a
solution. From left: a uniform stream $U_\infty$; a source of
strength $m$ with purely radial outflow $u_r = m/2\pi r$; a point
vortex of circulation $\Gamma$ with purely tangential flow
$u_\theta = \Gamma/2\pi r$; and the superposition of a uniform stream
with a doublet, which produces the closed streamline of a circular
cylinder. Adding a vortex to the last panel breaks the up-down
symmetry and yields the Kutta-Joukowski lift. The streamlines in the
rightmost panel are sketched from $\psi = U_\infty y(1-R^2/r^2)$.}
\label{fig:vol03ch11:potential-superposition}
\end{figure}
